\section{Experimental Setup}
\label{sec:setup}

\subsection{Tasks, Datasets and Protocols}

We evaluate on three tasks. \Cref{tab:datasets} lists the datasets and the
split used at each stage. Throughout, the \emph{selection} split is the one the
criterion $c_t$ is computed on, and it is always disjoint from the split used
for the reported numbers.

\begin{table*}[t]
  \centering
  \small
  \setlength{\tabcolsep}{5pt}
  \caption{Datasets and evaluation protocol. ``Selection'' is the split the
    ranking criterion $c_t$ is computed on; it is disjoint from ``Final
    evaluation'' in every case. Tiny ImageNet publishes no labelled test split,
    so its official validation split is the held-out evaluation set and its
    numbers are reported as validation metrics
    \cite{abai2019densenet,huynh2022vision}. For PASCAL VOC the original
    Ultralytics configuration aliases \texttt{val} to the VOC2007 test split,
    which would let the criterion see the evaluation data; we therefore hold out
  an independent selection split from the training pool.}
  \label{tab:datasets}
  \resizebox{\textwidth}{!}{%
\begin{tabular}{llccll}
    \toprule
    Task & Dataset & Classes & Available images & Selection & Final evaluation \\
    \midrule
    \multirow{5}{*}{Classification}
      & Tomato Leaf Disease \cite{solapure2024tomato}      & 10  & 1,609  & 15\% of set & 15\% test split \\
      & Burmese Grape Leaf Disease \cite{rahman2025burmese}& 5   & 3,103  & 15\% of set & 15\% test split \\
      & Potato Leaf Disease \cite{Shabrina2023Potato}      & 7   & 3,076  & 15\% of set & 15\% test split \\
      & CIFAR-100 \cite{krizhevsky2009learning}            & 100 & 50,000 train + 10,000 test & 5,000 from train & 10,000 official test \\
      & Tiny ImageNet \cite{le2015tiny}                    & 200 & 100,000 train + 10,000 val & 10,000 from train & 10,000 official val.\ \\
    \midrule
    Detection
      & PASCAL VOC \cite{everingham2010pascal}             & 20  & 16,551 trainval + 4,952 test & 1,655 from trainval & 4,952 VOC2007 test \\
    \midrule
    Segmentation
      & Carparts \cite{ultralytics2024carparts}            & 23  & 3,156 train + 401 val + 276 test & held out from train & 276 test split \\
    \bottomrule
  \end{tabular}}
\end{table*}

\paragraph{Classification.}
Three agricultural datasets---Tomato, Burmese Grape and Potato leaf
disease---provide a domain-specific, small-data regime, and CIFAR-100 and Tiny
ImageNet provide standard general-purpose benchmarks. The agricultural sets use
a stratified 70/15/15 train/selection/test partition. CIFAR-100 keeps its
official test split intact and carves 5,000 selection images out of the
training partition. Tiny ImageNet has no publicly labelled test split, so
following common practice on this benchmark
\cite{abai2019densenet,huynh2022vision} its official validation split serves as
the held-out evaluation set; every Tiny ImageNet number we report is therefore
labelled a validation metric, and its selection split is a separate 10,000
images stratified out of the training partition.

\paragraph{Detection.}
PASCAL VOC \cite{everingham2010pascal} in the Ultralytics configuration trains
on the union of the VOC2007 and VOC2012 \texttt{trainval} sets (16,551 images,
20 classes) and evaluates on the VOC2007 test set (4,952 images). That
configuration aliases the \texttt{val} split to the test split, which would
make the selection criterion a function of the evaluation data. We therefore
hold out 10\% of the training pool as an independent selection split and leave
the VOC2007 test set untouched for reporting.

\paragraph{Segmentation.}
Carparts \cite{ultralytics2024carparts} is an instance-segmentation dataset of
vehicle parts with 23 classes and a 3,156/401/276 train/val/test split. As
above, the selection split is held out from the training pool and the test
split is used only for reporting.

\subsection{Models}

For classification we fine-tune six ImageNet-pretrained backbones end-to-end,
instantiated from \texttt{timm} \cite{rw2019timm}, chosen to span
architectural families: VGG16 \cite{simonyan2015vgg} (plain convolutions),
ResNet101 \cite{he2016deep} (residual), DenseNet121 \cite{huang2017densely}
(dense connectivity), MobileNetV2 \cite{sandler2018mobilenetv2} (depthwise
separable, mobile budget), EfficientNetB0 \cite{tan2019efficientnet} (compound
scaling), and ViT-Base \cite{dosovitskiy2021imageworth16x16words}
(transformer). For detection and segmentation we use YOLOv8s
\cite{jocher2023ultralytics} and YOLO11s \cite{jocher2024yolo11}, covering two
generations of the same family. The six classification backbones across five
datasets give 30 configurations; two detectors on VOC and two segmentation
models on Carparts give four more, for 34 in total.

\subsection{Matched-Trajectory Protocol}

Every configuration is trained independently with five seeds
$\{1,10,42,100,500\}$. For a given seed, the baseline and all $k$ are derived
from the \emph{same} completed run and the same checkpoint set, so they share
architecture, initialization, data order, optimizer, schedule and training
duration, and differ only in the post-training selection rule. Any difference
between them is therefore attributable to that rule and not to training
variance. We report mean $\pm$ standard deviation over the five seeds. All
values of $k$ are reported as a sensitivity analysis; $k$ is never chosen using
the final evaluation split.

\subsection{Training Configuration}

Training pipelines are held fixed between the baseline and every $k$, so the
configuration below affects both arms identically and cannot favour either.
\Cref{tab:configs} records the per-model settings.

\begin{table*}[t]
  \centering \small
  \setlength{\tabcolsep}{5pt}\renewcommand{\arraystretch}{1.08}
  \caption{Per-model training configuration. The pipeline is identical for the baseline and every $k$, so these settings affect both arms equally and cannot favour either. \TODO{fill in}}
  \label{tab:configs}
  \resizebox{\textwidth}{!}{%
  \begin{tabular}{@{}llcccc@{\hspace{2.5em}}llcccc@{}}
    \toprule
    Dataset & Model & Batch & Epochs & LR & Patience & Dataset & Model & Batch & Epochs & LR & Patience \\
    \midrule
    \multirow{6}{*}{Tomato Leaf} & ViT-Base & & & & & \multirow{6}{*}{CIFAR-100} & ViT-Base & & & & \\
     & VGG16 & & & & &  & VGG16 & & & & \\
     & ResNet101 & & & & &  & ResNet101 & & & & \\
     & MobileNetV2 & & & & &  & MobileNetV2 & & & & \\
     & DenseNet121 & & & & &  & DenseNet121 & & & & \\
     & EfficientNetB0 & & & & &  & EfficientNetB0 & & & & \\\cmidrule(r){1-6}\cmidrule(l){7-12}
    \multirow{6}{*}{Burmese Grape} & ViT-Base & & & & & \multirow{6}{*}{Tiny ImageNet} & ViT-Base & & & & \\
     & VGG16 & & & & &  & VGG16 & & & & \\
     & ResNet101 & & & & &  & ResNet101 & & & & \\
     & MobileNetV2 & & & & &  & MobileNetV2 & & & & \\
     & DenseNet121 & & & & &  & DenseNet121 & & & & \\
     & EfficientNetB0 & & & & &  & EfficientNetB0 & & & & \\\cmidrule(r){1-6}\cmidrule(l){7-12}
    \multirow{6}{*}{Potato Leaf} & ViT-Base & & & & & \multirow{2}{*}{PASCAL VOC} & YOLOv8s & & & & \\
     & VGG16 & & & & &  & YOLO11s & & & & \\\cmidrule(l){7-12}
     & ResNet101 & & & & & \multirow{2}{*}{Carparts} & YOLOv8s & & & & \\
     & MobileNetV2 & & & & &  & YOLO11s & & & & \\
     & DenseNet121 & & & & &  &  & & & & \\
     & EfficientNetB0 & & & & &  &  & & & & \\\cmidrule(r){1-6}
    \bottomrule
  \end{tabular}}
\end{table*}

\paragraph{Shared settings.}
\TODO{fill in: optimizer, LR schedule and warmup, weight decay, loss,
input resolution, classifier head, augmentation, BN-recalibration batches.}
